\documentclass[12pt,a4paper]{amsart}
%\input xy
%\xyoption{all}
\usepackage[applemac]{inputenc}

\DeclareMathAlphabet{\mathpzc}{OT1}{pzc}{m}{it}
\usepackage{amsfonts,amsmath,amssymb,indentfirst,mathrsfs,amscd,url,breakurl}
\usepackage[mathscr]{eucal} 
%\usepackage[active]{srcltx}




\setlength{\oddsidemargin}{0.4in}
\setlength{\evensidemargin}{0.4in}
\setlength{\textwidth}{5.5in}
\setlength{\textheight}{8.8in}
\setlength{\marginparwidth}{0.8in}
\addtolength{\headheight}{2.5pt}

\renewcommand{\thefootnote}{}

\begin{document}
\thispagestyle{empty}
\vspace*{-2cm}

 
\noindent
\begin{center}\Large\bfseries\textsf{\'Algebra (grado en ingenier\'ia informatica)\\   curso 2013-14, primer semestre\\
	grupo 112}
\end{center}
\hrulefill{}
 
\noindent
{\sffamily    
Profesor:} Marco Zambon. \\
%{\sffamily Despacho:} 01.17.AU.103\\
{\sffamily Tutor\'ias:} por cita previa, en el despacho 605 del departamento de matem\'aticas (modulo 17). \\
{\sffamily P\'agina web del curso:} \url{http://www.uam.es/marco.zambon/WS13Ing.html}

\noindent\hrulefill{}  

\noindent
{\sffamily Horario:} Martes y Mi\'ercoles, 9:00-9:50, Viernes 11:00-12:50.\\
% Aula  XXX.   \\ 
\noindent
{\sffamily  Fecha de examenes:} \\ 
controles intermedios: habr\'a tres, en fechas a decidir.\\
%26 de octubre 2012 (viernes) y
% 30 de noviembre 2012 (viernes), en clase.\\ 
examen final: 13 de enero de 2014 (lunes, ma\~nana)\\
convocatoria extraordinaria:  10 de junio de 2014 (lunes, ma\~nana) 

\noindent\hrulefill{}  


\noindent
{\sffamily  Sistema de evaluaci\'on:}  (Ver la Guia Docente para m\'as detalles.)\\
%Hay dos opciones (itinerarios) para calcular la nota final.
%La manera m\'as conveniente me parece ser el itinerario 1, es decir, 
Hay dos opciones:
\begin{description}
\item[Itinerario 1] La nota se determina a partir del siguiente promedio:\\
controles intermedios  (12+12+16 por ciento),\\
examen final (60 por ciento).
\item[Itinerario 2] La nota es determinada por el examen final.
\end{description}

%Hay que presentarse a almenos 2 de las 4 pruebas.
%Los enunciados de los ejercicios realizados en clase consisteran de una selecci\'on de las hoja   para casa.

\noindent\hrulefill{}  

\noindent
{\sffamily  Hojas para casa:}\\
Habr\'a  hojas de ejercicios para casa, que discutiremos durante las horas de pr\'acticas. No ser\'a necesario entregar los ejercicios.
%La participaci\'on durante las horas de pr\'acticas puede contribuir a subir la nota hasta un 10 por ciento.

%\noindent
%{\sffamily  Tareas para casa:} Cada 2 semanas habr\'a una hoja de ejercicios para casa, que discutiremos en clase en la semana siguiente. No ser\'a necesario entregar los ejercicios.

%\noindent\hrulefill{}  


\noindent\hrulefill{}  

 \noindent
{\sffamily Programa orientativo del curso:} \\
Ver la Guia Docente.\\
Los bloques se presentar\'an en este orden: Bloque I, Bloque III, Bloque II.
  

\noindent\hrulefill{}  

 \noindent
{\sffamily Bibliograf\'ia:} \\
Ver la Guia Docente para una lista completa de libros aconsejados. La referencias principales ser\'an:


\begin{itemize}
\item \emph{F. Chamizo: \'Algebra I.} Notas de clase con listas de problemas.\\
Disponible en\\
 \url{http://www.uam.es/personal_pdi/ciencias/
 fchamizo/libreria/fich/APalgebraIinf96.pdf}

\item\emph{J. Dorronsoro y E. Hern\'andez: N�meros, grupos y anillos. Addison Wesley Iberoamericana, 1996}\\
(Cap�tulos  1 y 2, para los bloques I, III de la asignatura).

\item \emph{Sergei Treil. Linear algebra done wrong.}\\
 (Para el bloque II).\\ Disponible en\\
\url{http://www.math.brown.edu/~treil/papers/LADW/LADW.pdf}
 
  \end{itemize}

   
  

  
  
\end{document}
